% =============================================================================
% pthreads_learning_material.tex
% Learning Material: POSIX Threads (Pthreads)
% EC7207 High Performance Computing – Group 11
% =============================================================================
\documentclass[12pt,a4paper]{article}

\usepackage[margin=2.5cm]{geometry}
\usepackage{listings}
\usepackage{xcolor}
\usepackage{amsmath}
\usepackage{booktabs}
\usepackage{hyperref}
\usepackage{tikz}
\usepackage{array}
\usepackage{enumitem}
\usepackage{fancyhdr}
\usepackage{titlesec}

\lstdefinestyle{cstyle}{
    language=C,
    basicstyle=\ttfamily\small,
    keywordstyle=\color{blue}\bfseries,
    commentstyle=\color{gray}\itshape,
    stringstyle=\color{red!70!black},
    numbers=left, numberstyle=\tiny\color{gray},
    numbersep=8pt, frame=single, rulecolor=\color{black!30},
    breaklines=true, showstringspaces=false,
    backgroundcolor=\color{gray!5},
    tabsize=4
}
\lstset{style=cstyle}

\pagestyle{fancy}
\fancyhf{}
\lhead{EC7207 HPC – Group 11}
\rhead{POSIX Threads (Pthreads)}
\cfoot{\thepage}

\titleformat{\section}{\large\bfseries}{}{0em}{}[\titlerule]
\titleformat{\subsection}{\normalsize\bfseries}{}{0em}{}

% =============================================================================
\begin{document}
% =============================================================================

\begin{center}
    {\LARGE\bfseries POSIX Threads (Pthreads) — Manual Thread Management}\\[0.5em]
    {\large Learning Material for EC7207 High Performance Computing}\\[0.3em]
    {\normalsize Group 11 — Pearson Correlation Recommender System}
\end{center}
\vspace{1em}\hrule\vspace{1.5em}

% =============================================================================
\section{What are POSIX Threads?}
% =============================================================================

\textbf{Pthreads} (POSIX Threads) is the standard low-level threading API
for UNIX/Linux systems.  Unlike OpenMP, which uses compiler directives and
hides thread management, Pthreads requires the programmer to:

\begin{itemize}
    \item Explicitly \textbf{create} each thread with \texttt{pthread\_create()}
    \item Explicitly define what each thread does (a function with the signature
          \texttt{void* fn(void*)})
    \item Explicitly \textbf{divide} the work among threads
    \item Explicitly \textbf{wait} for threads to finish with
          \texttt{pthread\_join()}
\end{itemize}

\subsection{OpenMP vs.\ Pthreads — The Key Difference}

\begin{center}
\begin{tabular}{lll}
\toprule
Feature & OpenMP & Pthreads \\
\midrule
Thread creation & Automatic (\texttt{\#pragma}) & Manual (\texttt{pthread\_create}) \\
Work distribution & Automatic (\texttt{schedule}) & Manual (divide indices) \\
Synchronisation & Automatic (barriers) & Manual (\texttt{pthread\_join}) \\
Code changes & Minimal (add pragmas) & Significant (refactor into functions) \\
Portability & C/C++/Fortran standard & POSIX systems (Linux, macOS) \\
Control & Limited (clause-based) & Full (arbitrary logic) \\
\bottomrule
\end{tabular}
\end{center}

\textbf{Why learn Pthreads?}  OpenMP is convenient but opaque.  Pthreads
shows you \emph{exactly} what happens underneath: how threads are created,
how work is partitioned, and how synchronisation is achieved.  Understanding
Pthreads makes you a better parallel programmer even when using higher-level
APIs.

% =============================================================================
\section{Core Pthreads API}
% =============================================================================

\begin{center}
\begin{tabular}{lp{9cm}}
\toprule
Function & Purpose \\
\midrule
\texttt{pthread\_create(\&tid, attr, fn, arg)} &
  Create a new thread that runs \texttt{fn(arg)}.
  \texttt{attr = NULL} uses default attributes. \\
\texttt{pthread\_join(tid, \&retval)} &
  Wait for thread \texttt{tid} to finish.
  Returns the thread's return value via \texttt{retval}. \\
\texttt{pthread\_mutex\_lock(\&m)} &
  Acquire mutex \texttt{m} (blocks until available). \\
\texttt{pthread\_mutex\_unlock(\&m)} &
  Release mutex \texttt{m}. \\
\texttt{pthread\_barrier\_wait(\&b)} &
  Wait until all participating threads reach this point. \\
\bottomrule
\end{tabular}
\end{center}

Our project uses only \texttt{pthread\_create} and \texttt{pthread\_join}
because our work partitioning is disjoint — no thread ever writes to an index
another thread writes to, so no mutex is needed.

% =============================================================================
\section{Work Partitioning: \texttt{work\_range()}}
% =============================================================================

The key challenge in Pthreads is deciding \emph{which} thread does
\emph{which} work.  We use \textbf{ceiling division} to divide $N$ items
across $T$ threads:

\begin{lstlisting}
static void work_range(int total, int tid, int nthreads,
                       int *start, int *end)
{
    /* Ceiling division: (total + nthreads - 1) / nthreads */
    int chunk = (total + nthreads - 1) / nthreads;
    *start = tid * chunk;
    *end   = *start + chunk;
    if (*end > total) *end = total;  /* last thread may get fewer */
}
\end{lstlisting}

\textbf{Example} ($N = 10$ users, $T = 3$ threads):
\begin{center}
\begin{tabular}{clll}
\toprule
Thread & \texttt{start} & \texttt{end} & Users owned \\
\midrule
0 & 0 & 4 & 0, 1, 2, 3 \\
1 & 4 & 8 & 4, 5, 6, 7 \\
2 & 8 & 10 & 8, 9 (last thread, fewer) \\
\bottomrule
\end{tabular}
\end{center}

Ceiling division ensures all $N$ items are covered with no gaps and no
overlaps.  The last thread may get up to $T-1$ fewer items than the others —
this is the maximum load imbalance.

% =============================================================================
\section{Thread Argument Structure: \texttt{ThreadArgs}}
% =============================================================================

\texttt{pthread\_create} passes a single \texttt{void*} to the thread function.
We pack all the information a thread needs into a struct:

\begin{lstlisting}
typedef struct {
    int tid;       /* this thread's 0-based index */
    int nthreads;  /* total threads launched      */
} ThreadArgs;
\end{lstlisting}

From \texttt{tid} and \texttt{nthreads}, every thread can independently
compute its \texttt{[start, end)} work range.  No global variables are needed
for work distribution.

% =============================================================================
\section{Thread Functions: One Per Phase}
% =============================================================================

\subsection{Phase 2 — User Means Thread}

\begin{lstlisting}
static void *thread_user_means(void *arg)
{
    ThreadArgs *a = (ThreadArgs *)arg;
    int start, end;
    work_range(N_USERS, a->tid, a->nthreads, &start, &end);

    for (int u = start; u < end; u++) {
        double sum = 0.0; int cnt = 0;
        for (int i = 0; i < N_ITEMS; i++)
            if (R(u, i) != 0.0f) { sum += R(u, i); cnt++; }
        user_mean[u] = (cnt > 0) ? (float)(sum / cnt) : 3.0f;
    }
    return NULL;  /* thread function must return void* */
}
\end{lstlisting}

\textbf{Safety:}  Thread $t$ writes only to \texttt{user\_mean[start..end)}.
No two threads share an index.  No mutex needed.

\subsection{Phase 3 — Similarity Thread}

\begin{lstlisting}
static void *thread_similarities(void *arg)
{
    ThreadArgs *a = (ThreadArgs *)arg;
    int start, end;
    work_range(N_USERS, a->tid, a->nthreads, &start, &end);

    for (int u = start; u < end; u++) {
        SIM(u, u) = 1.0f;
        for (int v = u + 1; v < N_USERS; v++) {
            float s = pearson_similarity(u, v);
            SIM(u, v) = s;
            SIM(v, u) = s;  /* mirror -- race-free (same analysis as OpenMP) */
        }
    }
    return NULL;
}
\end{lstlisting}

\textbf{Race-free analysis:}  Thread $t$ owns outer loop values $u \in
[\text{start}, \text{end})$.  It writes \texttt{SIM(u,v)} and
\texttt{SIM(v,u)} for $v > u$.  No other thread processes the same outer
$u$ value.  When another thread (say thread $s$ owning $u' > u$) writes
\texttt{SIM(u',v')} for $v' > u'$, it never touches the cell
\texttt{SIM(v,u)} because $v > u$ and $u' > u$ means $u' \neq u$.

\subsection{Phase 4 — Prediction Thread}

\begin{lstlisting}
static void *thread_predictions(void *arg)
{
    ThreadArgs *a = (ThreadArgs *)arg;
    int start, end;
    work_range(N_USERS, a->tid, a->nthreads, &start, &end);

    /* Private scratch buffer for this thread only */
    SimPair *nbrs = (SimPair *)malloc(N_USERS * sizeof(SimPair));

    for (int u = start; u < end; u++) {
        for (int item = 0; item < N_ITEMS; item++) {
            /* ... prediction logic using nbrs ... */
        }
    }

    free(nbrs);
    return NULL;
}
\end{lstlisting}

Each thread allocates its own \texttt{nbrs} buffer on the heap.  Unlike a
global buffer, this is never shared between threads.

% =============================================================================
\section{The \texttt{run\_parallel()} Harness}
% =============================================================================

A helper function handles the create/join pattern for each phase:

\begin{lstlisting}
static double run_parallel(void *(*fn)(void *),
                           pthread_t  *threads,
                           ThreadArgs *args)
{
    double t0 = now_sec();

    /* Create all N_THREADS threads */
    for (int t = 0; t < N_THREADS; t++) {
        args[t].tid      = t;
        args[t].nthreads = N_THREADS;
        pthread_create(&threads[t], NULL, fn, &args[t]);
    }

    /* Wait for all threads to complete */
    for (int t = 0; t < N_THREADS; t++)
        pthread_join(threads[t], NULL);

    return now_sec() - t0;  /* wall-clock time for this phase */
}
\end{lstlisting}

\textbf{The join is a barrier:}  \texttt{pthread\_join} returns only after
the named thread has exited.  Joining all $T$ threads sequentially guarantees
that Phase $n+1$ starts only after every thread of Phase $n$ has finished.
This is equivalent to the implicit barrier in OpenMP's \texttt{parallel for}.

% =============================================================================
\section{Phase Control Flow in \texttt{main()}}
% =============================================================================

\begin{lstlisting}
pthread_t  threads[MAX_THREADS];
ThreadArgs args[MAX_THREADS];

/* Stack-allocated: no malloc needed for these arrays */
generate_data();                            /* serial */
run_parallel(thread_user_means,  threads, args);   /* parallel */
run_parallel(thread_similarities, threads, args);  /* parallel */
run_parallel(thread_predictions,  threads, args);  /* parallel */
evaluate_mae();                             /* serial */
\end{lstlisting}

Threads are created fresh for each phase.  This re-creation overhead
(\texttt{pthread\_create} takes $\sim$10--100 \textmu s) is negligible
compared to the similarity computation ($\sim$seconds).

An alternative is a \textbf{thread pool} with condition variables (threads
sleep between phases and wake on a signal).  We chose the simpler
create/join pattern because it makes phase ordering obvious and the overhead
is trivially small.

% =============================================================================
\section{Data Races — What They Are and How We Avoid Them}
% =============================================================================

\subsection{What is a Data Race?}

A data race occurs when:
\begin{enumerate}
    \item Two or more threads access the same memory location concurrently,
    \item At least one access is a write, and
    \item There is no synchronisation (mutex, atomic, barrier) ordering the accesses.
\end{enumerate}

Data races produce \textbf{undefined behaviour} — the result may be wrong,
different on every run, or apparently correct but broken on a different CPU.

\subsection{How Our Code Avoids Data Races}

\begin{center}
\begin{tabular}{lll}
\toprule
Array & Access & Race-free because \\
\midrule
\texttt{ratings[]} & Read only after Phase 1 & Multiple readers are safe \\
\texttt{user\_mean[]} & Written per-thread & Disjoint index ranges \\
\texttt{sim\_matrix[]} & Written per-thread & Disjoint outer-loop rows \\
\texttt{predictions[]} & Written per-thread & Disjoint row ranges \\
\texttt{nbrs[]} & Per-thread private allocation & Never shared \\
\bottomrule
\end{tabular}
\end{center}

\subsection{If We Had a Bug — What Would Happen?}

If two threads wrote to the same \texttt{sim\_matrix} cell, the result would
be one of the two values (whichever write ``won'' the race).  The
sim-matrix checksum would then differ from the serial baseline.  Our
\texttt{[Check]} output line catches this.

% =============================================================================
\section{Pthreads vs.\ OpenMP for This Project}
% =============================================================================

\begin{center}
\begin{tabular}{lll}
\toprule
Aspect & Pthreads & OpenMP \\
\midrule
Lines of new code & $\sim$50 extra (run\_parallel, thread fns) & $\sim$6 extra pragmas \\
Scheduling & Manual (static only as implemented) & static / dynamic / guided \\
Load balancing & Our \texttt{work\_range()} divides equally & Dynamic handles unequal work \\
Barrier & Explicit (\texttt{pthread\_join} loop) & Implicit (end of parallel for) \\
Debugging & Easier (see every thread creation) & Harder (abstracted away) \\
\bottomrule
\end{tabular}
\end{center}

\textbf{Why does OpenMP potentially outperform our Pthreads for Phase 3?}
OpenMP's \texttt{schedule(dynamic)} adapts to the decreasing inner-loop
size.  Our Pthreads version uses static chunking of outer rows, which
can leave the thread assigned to small-$u$ rows (more work) as the
bottleneck while others sit idle.

% =============================================================================
\section{Compile and Run}
% =============================================================================

\begin{verbatim}
# Compile (-lpthread links the pthreads library)
gcc -O2 -Wall -o pthreads_rec pthreads/pthreads_recommender.c -lpthread -lm

# Run with different thread counts
./pthreads_rec 1000 1000 1    # 1 thread (serial-equivalent)
./pthreads_rec 1000 1000 2    # 2 threads
./pthreads_rec 1000 1000 4    # 4 threads
./pthreads_rec 1000 1000 8    # 8 threads
\end{verbatim}

Arguments: \texttt{num\_users num\_items num\_threads}

% =============================================================================
\section{Key Points to Explain in a Presentation}
% =============================================================================

\begin{enumerate}
    \item \textbf{Why create threads fresh per phase, not once at the start?}
          Simplicity and clarity.  Creating threads once would require condition
          variables to signal work; the overhead of create/join for
          3 phases is $\ll 1$\% of compute time.
    \item \textbf{What is the \texttt{void* fn(void*)} signature?}  Pthreads
          requires a uniform function signature for the thread entry point.  The
          \texttt{void*} argument is a generic pointer that we cast to
          \texttt{ThreadArgs*} inside the function.
    \item \textbf{Why is \texttt{pearson\_similarity()} safe to call from multiple threads?}
          It reads only from \texttt{ratings[]} and \texttt{user\_mean[]},
          which are fully written before the parallel phase begins, and writes
          only to local variables (stack-allocated scalars).  This makes it a
          pure function with no side effects — inherently thread-safe.
    \item \textbf{How does ceiling division handle non-divisible sizes?}
          For $N=1000, T=3$: chunk $= \lceil 1000/3 \rceil = 334$.
          Thread 0: 0--333, Thread 1: 334--667, Thread 2: 668--999 (332 items).
          All covered, no gaps, no overlaps.
\end{enumerate}

\vfill
\begin{center}
    \small\textit{Group 11 — EC7207 High Performance Computing}
\end{center}

\end{document}
